% Chapter 1 -- The bootstrap.  Sources (Law 1): ADR-0007 DEC-01 (two
% compilers, never conflated), DEC-06 (Phi, K1/K2/K3, the K2 == K3
% certificate, scope, reproducibility precondition, DDC), DEC-08 (M7 is
% seed-built; resident self-compilation is M8); PRD Sec 7; CLAUDE.md
% hallucination-risk callouts (K2 == K3, not K1 == K2).
\Gates{test-tps-lex-build, test-tps-gen-fpc, test-compiler, test-pubs-tps}
\chapter{The bootstrap}
\ix{bootstrap}

\section{Two compilers}
\ix{genesis compiler}
\ix{resident compiler}
Two compilers take part in building Turbo Initech, and they are never
confused with each other:

\SideHead{The genesis compiler}\file{INITECHC}, the genesis compiler, is a
Pascal compiler written in C that runs on the development system.  It
compiles the Turbo Initech source once, to make the first Turbo Initech
compiler.  Its only purpose is to start the chain; an error in it is an
error of the genesis compiler, not of Turbo Initech, and it is not shipped.

\SideHead{Turbo Initech}Turbo Initech, the resident compiler, is written in
Pascal.  It compiles itself and the Pascal programs written for InitechOS.

The InitechDOS operating system and the other parts of InitechOS are not
written in Pascal and are not built by either compiler; they are built by
the development system's own tools.  Nothing in this book applies to them.

\section{Compiling a compiler}
\ix{self-compilation}
Write \textit{src} for the Turbo Initech source, and $K(\textit{src})$ for
the program that compiler $K$ makes from it.  The bootstrap is the chain:

\begin{BorTable}{The chain}{p{24mm}p{74mm}}
\Hd{Step} & \Hd{Made by}\EndHd
$K_1 = X(\textit{src})$ & the genesis compiler $X$ compiles the source: the
  first Turbo Initech compiler\\
$K_2 = K_1(\textit{src})$ & the first Turbo Initech compiler compiles its
  own source\\
$K_3 = K_2(\textit{src})$ & the compiler so made compiles the source
  again\\
\end{BorTable}

\ix{fixed point}
\ix{certificate}
Each Turbo Initech compiler in the chain compiles the same source with the
same rules, so from $K_2$ on the chain must repeat itself: $K_3$ must equal
$K_2$, byte for byte.  This equality is the \term{self-compilation
certificate}.

\ArrowNote{The certificate is $K_2 = K_3$, not $K_1 = K_2$.  $K_1$ is made by
the genesis compiler, a different program, and may differ from $K_2$ in
every byte without anything being wrong.}

The certificate means something only when every build is
\term{reproducible}: the same compiler and the same source must always
give the same bytes, with no date, time, directory or chance in the output
(Chapter 6).

\section{Independent compilation}
\ix{diverse double-compilation}
A compiler that compiles itself could, in principle, carry a fault from
one generation to the next by reproducing it on purpose.  The defense is
to make the first compiler a second time with a genesis compiler written
independently, and to check that the chains of the two meet (Chapter 5).

\section{Where the chain stands}
\ix{status}
The first steps of the chain are built and checked today.  The genesis
compiler makes $K_1$, the Turbo Initech compiler \file{TPS.COM}, and $K_1$
runs on InitechDOS.  On the development system, a build of Turbo Initech
made by the Free Pascal compiler for testing compiles the Turbo Initech
source, twice with the same bytes, into assembly language that links into a
complete program of the right size.
\draft{the certificate K2 == K3 is not yet produced: make selfhost is an M8
placeholder (ADR-0007 DEC-08, bead initech-ls4); it needs the resident
assembler and linker}
Chapter 4 says, step by step, which parts of the acceptance procedure are
carried out today.
